\documentclass[12pt,a4paper]{article}
% Preámbulo del informe de la Etapa 2.
%
% Tipografía: Computer Modern a 12pt, que es exactamente la que usa la cátedra
% en «Reglas de Sintaxis de MiniJava» (pdffonts sobre ese PDF reporta CMR12 para
% el cuerpo, CMCSC10 para las versalitas, CMTT10 para las producciones y
% CMSSBX10 para los terminales). `lmodern` es Latin Modern, la reimplementación
% Unicode del mismo diseño.
%
% Paleta monocromática y estructura heredadas del informe de la Etapa 1.
\usepackage[T1]{fontenc}
\usepackage[utf8]{inputenc}
\usepackage{lmodern}
\usepackage{amsmath,amssymb}
\usepackage[spanish,es-noquoting,es-noshorthands,es-tabla]{babel}
\usepackage[a4paper,margin=2.6cm,headsep=14pt]{geometry}
\usepackage{microtype}
\usepackage{booktabs}
\usepackage{longtable}
\usepackage{array}
\usepackage{xcolor}
\usepackage{fancyhdr}
\usepackage{titlesec}
\usepackage{enumitem}
\usepackage{fancyvrb}
\usepackage{framed}
\usepackage[hidelinks,bookmarks=true]{hyperref}

% Caracteres que aparecen en los comentarios de los archivos de gramática y
% llegan al informe a través de los listados generados.
\DeclareUnicodeCharacter{03B5}{\ensuremath{\epsilon}}
\DeclareUnicodeCharacter{2014}{---}
\DeclareUnicodeCharacter{2026}{\ldots}

% --- Paleta monocromática: sólo grises, con un único acento. -----------------
\definecolor{tinta}{gray}{0.13}
\definecolor{grisMedio}{gray}{0.42}
\definecolor{grisSuave}{gray}{0.70}
\definecolor{grisFondo}{gray}{0.955}
\definecolor{acento}{HTML}{2F4858}

\color{tinta}

% --- Títulos ----------------------------------------------------------------
\titleformat{\section}{\Large\bfseries\color{acento}}{\thesection}{0.7em}{}
\titleformat{\subsection}{\large\bfseries\color{tinta}}{\thesubsection}{0.7em}{}
\titleformat{\subsubsection}{\bfseries\color{grisMedio}}{\thesubsubsection}{0.7em}{}
\titlespacing*{\section}{0pt}{2.4ex plus 1ex minus .2ex}{1.3ex plus .2ex}

% --- Encabezado y pie -------------------------------------------------------
\setlength{\headheight}{14pt}
\pagestyle{fancy}
\fancyhf{}
\renewcommand{\headrulewidth}{0.4pt}
\renewcommand{\headrule}{\hbox to\headwidth{\color{grisSuave}\leaders\hrule height \headrulewidth\hfill}}
\fancyhead[L]{\footnotesize\color{grisMedio}Compiladores e Intérpretes — Etapa 2}
\fancyhead[R]{\footnotesize\color{grisMedio}\nouppercase{\leftmark}}
\fancyfoot[C]{\footnotesize\color{grisMedio}\thepage}

% ===========================================================================
% Notación de la gramática — la misma convención tipográfica de la cátedra
% ===========================================================================
% \nt{X}   no terminal   -> \texttt, como <X> en el PDF de la cátedra
% \tl{x}   terminal      -> sans-serif negrita, como los terminales del PDF
% \eps     la cadena vacía
% \deriva  el ::= de las producciones
\newcommand{\nt}[1]{\texttt{<#1>}}
\newcommand{\tl}[1]{\textsf{\bfseries #1}}
\newcommand{\eps}{\ensuremath{\epsilon}}
\newcommand{\deriva}{\texttt{::=}}
\newcommand{\alt}{\ensuremath{\;\mid\;}}
\newcommand{\prim}[1]{\textup{Primeros}(#1)}

% Bloque de producciones: sangrado, interlineado propio, sin justificar.
\newenvironment{producciones}{%
  \par\smallskip\begingroup
  \parindent=0pt \leftskip=1.2em
  % Elasticidad acotada: evita estirar los espacios de una producción larga,
  % pero no tanta como para que el \hfill de la etiqueta de logro no empuje.
  \rightskip=0pt plus 2em
  \parfillskip=0pt plus 1fil
  \setlength{\parskip}{0.35ex}%
  \obeylines}{\par\endgroup\smallskip}

% Una producción con su nombre a la izquierda, para las tablas de Primeros.
\newcommand{\regla}[2]{\nt{#1} \deriva{} #2}

% --- Bloques de código y de consola -----------------------------------------
\DefineVerbatimEnvironment{consola}{Verbatim}%
  {fontsize=\small,xleftmargin=1.1em,formatcom=\color{tinta}}
\DefineVerbatimEnvironment{codigo}{Verbatim}%
  {fontsize=\small,xleftmargin=1.1em,commandchars=\\\{\}}

% --- Cajas de defecto -------------------------------------------------------
% Dos clases de defecto que NO hay que confundir; van con marcas distintas.
\newenvironment{conflicto}[1]{%
  \par\smallskip\begingroup
  \setlength{\fboxsep}{0.7em}%
  \definecolor{shadecolor}{gray}{0.955}%
  \begin{shaded}\noindent{\bfseries\color{acento}Conflicto LL(1) --- #1}\par\smallskip
}{\end{shaded}\endgroup\smallskip}

\newenvironment{ambiguedad}[1]{%
  \par\smallskip\begingroup
  \setlength{\fboxsep}{0.7em}%
  \definecolor{shadecolor}{gray}{0.90}%
  \begin{shaded}\noindent{\bfseries\color{tinta}Ambigüedad --- #1}\par\smallskip
}{\end{shaded}\endgroup\smallskip}

\newenvironment{decision}[1]{%
  \par\smallskip\begingroup
  \setlength{\fboxsep}{0.7em}%
  \definecolor{shadecolor}{gray}{0.975}%
  \begin{shaded}\noindent{\bfseries\color{grisMedio}Decisión --- #1}\par\smallskip
}{\end{shaded}\endgroup\smallskip}

% --- Los cuatro pasos de cada transformación --------------------------------
\newcommand{\paso}[1]{\par\smallskip\noindent\textbf{\color{grisMedio}#1}\quad\ignorespaces}

% Los nombres largos de no terminales van en \texttt y no se parten, así que se
% le da al justificado un margen de maniobra en vez de dejar renglones desbordados.
\setlength{\emergencystretch}{2.5em}

\setlist[itemize]{leftmargin=1.3em,itemsep=0.25ex,topsep=0.5ex}
\setlist[enumerate]{leftmargin=1.5em,itemsep=0.25ex,topsep=0.5ex}
\setlist[description]{leftmargin=1.3em,itemsep=0.35ex,topsep=0.6ex}
\renewcommand{\arraystretch}{1.25}

\usepackage{ragged2e}
\usepackage{needspace}

% Conjuntos de Primeros que se repiten en el paso a paso.
\newcommand{\cT}{\ensuremath{\mathcal{T}}}
\newcommand{\cE}{\ensuremath{\mathcal{E}}}
\newcommand{\cP}{\ensuremath{\mathcal{P}}}
\newcommand{\cj}[1]{\ensuremath{\{}#1\ensuremath{\}}}
\newcommand{\tec}[1]{\textsf{\bfseries #1}}
\newcommand{\ok}{\ensuremath{\checkmark}}

\begin{document}

% ===========================================================================
% Portada
% ===========================================================================
\begin{titlepage}
\thispagestyle{empty}
\raggedright
\vspace*{2.2cm}

{\footnotesize\sffamily
UNIVERSIDAD NACIONAL DEL SUR\\
Departamento de Ciencias e Ingeniería de la Computación\\
Compiladores e Intérpretes --- Segundo Cuatrimestre de 2026\par}

\vspace{2.4cm}
{\rule{\textwidth}{0.6pt}}
\vspace{1.1cm}

{\Huge\bfseries Compilador de MiniJava\par}
\vspace{0.5cm}
{\LARGE\color{tinta} Etapa 2 --- Analizador Sintáctico\\[0.25cm]
\large Descendente Recursivo Predictivo\par}

\vspace{1.1cm}
{\rule{\textwidth}{0.6pt}}

\vspace{2.6cm}
{\large Agustín Paul\par}
\vspace{0.2cm}
{\normalsize L.U. 147524\par}
\vspace{0.4cm}
{Informe de entrega\par}

\vfill
{\color{tinta}\footnotesize
Transformación de la gramática a LL(1), paso a paso.\par}
\vspace{0.6cm}
{\color{tinta}\footnotesize \today\par}
\end{titlepage}

\thispagestyle{empty}
\tableofcontents
\newpage
\setcounter{page}{1}
\RaggedRight

% ===========================================================================
\section{Cómo compilar y ejecutar}
\label{sec:uso}
% ===========================================================================

Proyecto Maven con Java~25. Desde la carpeta \texttt{compilador/}:

\begin{consola}
$ mvn package          # compila, corre la suite y arma target/Compilador.jar
$ java -jar target/Compilador.jar programa1.java
Compilacion Exitosa

[SinErrores]
\end{consola}

El fuente es el único argumento y puede tener cualquier extensión. \texttt{-{}-etapa=1} corre sólo el
analizador léxico de la Etapa 1. El compilador corre siempre con todos sus logros.

% ===========================================================================
\section{Logros}
\label{sec:logros-perseguidos}
% ===========================================================================

\begin{center}
\begin{tabular}{@{}lp{0.60\linewidth}@{}}
\toprule
\textbf{Logro} & \textbf{Dónde está} \\
\midrule
Variables locales Clásicas E2 & L4 (§\ref{sec:l4}). \\
Visibilidad Mejorada E2 & L6 (§\ref{sec:l6}). Agrega el token \tl{private}. \\
Fors! E2 & L7 (§\ref{sec:l7}). Agrega el token \tl{for}. \\
Operador Ternario E2 & L3 (§\ref{sec:l3}). Agrega el token \tl{?}. \\
Operadores Posifijos E2 & L2 (§\ref{sec:l2}). \\
Atributos Inicializados & L5 (§\ref{sec:l5}). \\
Inicializacion de Arreglos E2 & L8 (§\ref{sec:l8}). \\
Genericidad Avanzada E2 & L9 (§\ref{sec:l9}). \\
Recuperación de Errores Modo Pánico & Implementación (§\ref{sec:errores}). \\
Floats I (Etapa 1) & L1 (§\ref{sec:l1}). \\
\bottomrule
\end{tabular}
\end{center}

También se persiguen \emph{Entrega Anticipada}, \emph{Imbatibilidad} y \emph{Well Done!}, que no
dependen del código.

% ===========================================================================
\section{Transformación de la gramática}
\label{sec:transformaciones}
% ===========================================================================

\subsection{Método}

\paso{Criterio LL(1).} Para todo $A \to \beta_1 \mid \dots \mid \beta_n$,
$\prim{\beta_i} \cap \prim{\beta_j} = \varnothing$ si $i \neq j$.

\paso{Orden.} Primero se elimina la recursión a izquierda y después se factoriza, porque eliminar la
recursión puede dejar producciones sin factorizar (R3).

\needspace{10\baselineskip}
\paso{Técnicas.}

\begin{center}\small
\begin{tabular}{@{}l>{\raggedright\arraybackslash}p{0.84\linewidth}@{}}
\toprule
\tec{RI} & Eliminación de recursión a izquierda directa:
$A \to A\alpha_1 \mid \dots \mid A\alpha_n \mid \beta_1 \mid \dots \mid \beta_m$ pasa a
$A \to \beta_1 R \mid \dots \mid \beta_m R$ y $R \to \alpha_1 R \mid \dots \mid \alpha_n R \mid \epsilon$. \\
\tec{FD} & Factorización a izquierda directa:
$X \to \alpha\beta_1 \mid \dots \mid \alpha\beta_n \mid \gamma_1 \mid \dots \mid \gamma_m$ pasa a
$X \to \alpha R \mid \gamma_1 \mid \dots \mid \gamma_m$ y $R \to \beta_1 \mid \dots \mid \beta_n$. \\
\tec{SU} & Subir: reemplazar un no terminal por sus partes derechas para exponer el prefijo común. \\
\tec{FP} & Factorización profunda: \tec{SU} hasta que el prefijo sea directo, y después \tec{FD}. \\
\tec{AG} & Agrupar: reemplazar alternativas, o una cola repetida, por un no terminal nuevo que deriva
exactamente eso. Es la inversa de \tec{SU}: no cambia el lenguaje ni hace falta para LL(1). Se usa
para que cada punto de decisión conserve un nombre, porque las reglas tienen un significado que el
semántico necesita. \\
\bottomrule
\end{tabular}
\end{center}

\paso{Conjuntos que se repiten.}
\cT{} $=\prim{\nt{Tipo}}=$ \cj{\tl{boolean}, \tl{char}, \tl{int}, \tl{idClase}, \tl{idGen}};
\cE{} $=\prim{\nt{Expresion}}=$ \cj{\tl{!}, \tl{(}, \tl{+}, \tl{-}, \tl{charLiteral}, \tl{false},
\tl{idClase}, \tl{idMetVar}, \tl{intLiteral}, \tl{new}, \tl{null}, \tl{stringLiteral}, \tl{this},
\tl{true}};
\cP{} $=\prim{\nt{Primario}}=$ \cj{\tl{(}, \tl{idClase}, \tl{idMetVar}, \tl{new}, \tl{stringLiteral},
\tl{this}}. Desde L1, \cT{} incluye \tl{float} y \cE{}, \tl{floatLiteral}.

\subsection{Diagnóstico}

La gramática de la cátedra tiene \textbf{16 conflictos en 10 no terminales}.

\begin{center}\small
\begin{tabular}{@{}c>{\raggedright\arraybackslash}p{0.28\linewidth}>{\raggedright\arraybackslash}p{0.33\linewidth}l l@{}}
\toprule
\textbf{\#} & \textbf{No terminal} & \textbf{Defecto} & \textbf{Intersección} & \textbf{Paso} \\
\midrule
1     & \nt{Miembro}               & prefijo común no directo & \cT & F6 \\
2     & \nt{ListaArgsFormales}     & recursión a izquierda    & \cT & R1 \\
3     & \nt{Sentencia}             & \nt{Asignacion} y \nt{Llamada} idénticos & \cE & F2 \\
4     & \nt{If}                    & prefijo común            & \cj{\tl{if}} & F3 \\
5     & \nt{Expresion}             & prefijo común            & \cE & F1 \\
6     & \nt{ExpresionCompuesta}    & recursión a izquierda    & \cE & R2 \\
7--12 & \nt{Referencia}            & recursión a izquierda, seis pares & \cP & R3 \\
13    & \nt{Primario}              & prefijo común no directo & \cj{\tl{idMetVar}} & F7 \\
14    & \nt{Primario}              & prefijo común a dos niveles & \cj{\tl{new}} & F8 \\
15    & \nt{DimensionesConTamanio} & prefijo común            & \cj{\tl{[}} & F4 \\
16    & \nt{ListaExps}             & prefijo común            & \cE & F5 \\
\bottomrule
\end{tabular}
\end{center}

\paso{Ambigüedades.} Un conflicto LL(1) no implica ambigüedad, que es tener dos árboles para una
misma cadena. La base tiene tres: (1) \nt{Asignacion} y \nt{Llamada} tienen el mismo cuerpo, así que
toda sentencia-expresión tiene dos árboles; se elimina en F2. (2) El \emph{dangling else},
\texttt{if (a) if (b) s1; else s2;}; sobrevive a F3. (3) En \texttt{new int[3][4]} el segundo
\tl{[} puede ser otra dimensión o un acceso al arreglo creado; no es un conflicto de Primeros de la
base y aparece en F4. En los otros catorce conflictos las alternativas derivan lenguajes disjuntos.

% ---------------------------------------------------------------------------
\subsection{Eliminación de recursión a izquierda}

\subsubsection*{R1 --- \nt{ListaArgsFormales}}

\paso{Original}
\begin{producciones}
\regla{ListaArgsFormales}{\nt{ArgFormal} \alt{} \nt{ListaArgsFormales} \tl{,} \nt{ArgFormal}}
\end{producciones}
\paso{Conflicto.} Las dos alternativas tienen Primeros \cT.

\paso{Paso.} \tec{RI} con $\beta_1 =$ \nt{ArgFormal}, $\alpha_1 =$ \tl{,} \nt{ArgFormal} y
$R =$ \nt{RestoListaArgsFormales}.

\paso{Resultado}
\begin{producciones}
\regla{ListaArgsFormales}{\nt{ArgFormal} \nt{RestoListaArgsFormales}}
\regla{RestoListaArgsFormales}{\tl{,} \nt{ArgFormal} \nt{RestoListaArgsFormales} \alt{} \eps}
\end{producciones}
\cj{\tl{,}} contra \cj{\eps}. \ok

\subsubsection*{R2 --- \nt{ExpresionCompuesta}}

\paso{Original}
\begin{producciones}
\regla{ExpresionCompuesta}{\nt{ExpresionCompuesta} \nt{OperadorBinario} \nt{ExpresionBasica} \alt{} \nt{ExpresionBasica}}
\end{producciones}
\paso{Conflicto.} Las dos alternativas tienen Primeros \cE.

\paso{Paso.} \tec{RI} con $\beta_1 =$ \nt{ExpresionBasica}, $\alpha_1 =$ \nt{OperadorBinario}
\nt{ExpresionBasica} y $R =$ \nt{RestoExpresionCompuesta}.

\paso{Resultado}
\begin{producciones}
\regla{ExpresionCompuesta}{\nt{ExpresionBasica} \nt{RestoExpresionCompuesta}}
\regla{RestoExpresionCompuesta}{\nt{OperadorBinario} \nt{ExpresionBasica} \nt{RestoExpresionCompuesta} \alt{} \eps}
\end{producciones}
Los trece operadores binarios contra \cj{\eps}. \ok{} Siguen en un solo nivel, como en la base.

\subsubsection*{R3 --- \nt{Referencia}}

\paso{Original}
\begin{producciones}
\regla{Referencia}{\nt{Primario} \alt{} \nt{Referencia} \nt{VarEncadenada} \alt{} \nt{Referencia} \nt{MetodoEncadenado} \alt{} \nt{Referencia} \nt{AccesoArreglo}}
\regla{VarEncadenada}{\tl{.} \tl{idMetVar}}
\regla{MetodoEncadenado}{\tl{.} \tl{idMetVar} \nt{ArgsActuales}}
\end{producciones}
\paso{Conflicto.} Las cuatro alternativas tienen Primeros \cP: seis pares.

\paso{Paso 1.} \tec{RI} con $\beta_1 =$ \nt{Primario}, $\alpha_1 =$ \nt{VarEncadenada},
$\alpha_2 =$ \nt{MetodoEncadenado}, $\alpha_3 =$ \nt{AccesoArreglo} y $R =$ \nt{RestoReferencia}:
\begin{producciones}
\regla{Referencia}{\nt{Primario} \nt{RestoReferencia}}
\regla{RestoReferencia}{\nt{VarEncadenada} \nt{RestoReferencia} \alt{} \nt{MetodoEncadenado} \nt{RestoReferencia} \alt{} \nt{AccesoArreglo} \nt{RestoReferencia} \alt{} \eps}
\end{producciones}
Conflicto nuevo: las dos primeras alternativas tienen Primeros \cj{\tl{.}}.

\paso{Paso 2 (\tec{AG}).} Las dos comparten la cola \nt{RestoReferencia}:
\begin{producciones}
\regla{RestoReferencia}{\nt{Encadenado} \nt{RestoReferencia} \alt{} \nt{AccesoArreglo} \nt{RestoReferencia} \alt{} \eps}
\regla{Encadenado}{\nt{VarEncadenada} \alt{} \nt{MetodoEncadenado}}
\end{producciones}
El conflicto pasa a \nt{Encadenado}: \cj{\tl{.}} contra \cj{\tl{.}}.

\paso{Paso 3.} \tec{FP}: \tec{SU} de \nt{VarEncadenada} y \nt{MetodoEncadenado},
\begin{producciones}
\regla{Encadenado}{\tl{.} \tl{idMetVar} \alt{} \tl{.} \tl{idMetVar} \nt{ArgsActuales}}
\end{producciones}
y \tec{FD} con $\alpha =$ \tl{.} \tl{idMetVar}, $\beta_1 = \epsilon$, $\beta_2 =$ \nt{ArgsActuales} y
$R =$ \nt{EncadenadoAtributoOMetodo}.

\paso{Resultado}
\begin{producciones}
\regla{Referencia}{\nt{Primario} \nt{RestoReferencia}}
\regla{RestoReferencia}{\nt{Encadenado} \nt{RestoReferencia} \alt{} \nt{AccesoArreglo} \nt{RestoReferencia} \alt{} \eps}
\regla{Encadenado}{\tl{.} \tl{idMetVar} \nt{EncadenadoAtributoOMetodo}}
\regla{EncadenadoAtributoOMetodo}{\nt{ArgsActuales} \alt{} \eps}
\end{producciones}
\nt{RestoReferencia}: \cj{\tl{.}}, \cj{\tl{[}} y \cj{\eps}. \nt{EncadenadoAtributoOMetodo}:
\cj{\tl{(}} contra \cj{\eps}. \ok{} Se pierden \nt{VarEncadenada} y \nt{MetodoEncadenado}; la
distinción queda en \nt{EncadenadoAtributoOMetodo}: vacío es el atributo y con argumentos, el método.

% ---------------------------------------------------------------------------
\subsection{Factorización}

\subsubsection*{F1 --- \nt{Expresion}}

\paso{Original}
\begin{producciones}
\regla{Expresion}{\nt{ExpresionCompuesta} \nt{OperadorAsignacion} \nt{ExpresionCompuesta} \alt{} \nt{ExpresionCompuesta}}
\end{producciones}
\paso{Conflicto.} Las dos alternativas tienen Primeros \cE.

\paso{Paso.} \tec{FD} con $\alpha =$ \nt{ExpresionCompuesta}, $\beta_1 =$ \nt{OperadorAsignacion}
\nt{ExpresionCompuesta}, $\beta_2 = \epsilon$ y $R =$ \nt{RestoExpresion}.

\paso{Resultado}
\begin{producciones}
\regla{Expresion}{\nt{ExpresionCompuesta} \nt{RestoExpresion}}
\regla{RestoExpresion}{\nt{OperadorAsignacion} \nt{ExpresionCompuesta} \alt{} \eps}
\end{producciones}
\cj{\tl{=}} contra \cj{\eps}. \ok{} El lado derecho del \tl{=} sigue siendo una
\nt{ExpresionCompuesta}: \texttt{a = b = c} no pertenece al lenguaje, como en la base.

\subsubsection*{F2 --- \nt{Sentencia}}

\paso{Original}
\begin{producciones}
\regla{Sentencia}{\tl{;} \alt{} \nt{Asignacion} \tl{;} \alt{} \nt{Llamada} \tl{;} \alt{} \nt{VarLocal} \tl{;} \alt{} \nt{Return} \tl{;} \alt{} \nt{If} \alt{} \nt{While} \alt{} \nt{Bloque}}
\regla{Asignacion}{\nt{Expresion}}
\regla{Llamada}{\nt{Expresion}}
\end{producciones}
\paso{Conflicto.} \nt{Asignacion} \tl{;} y \nt{Llamada} \tl{;} tienen Primeros \cE{}, y además son la
misma producción: la ambigüedad (1).

\paso{Paso 1.} \tec{SU} de \nt{Asignacion} y \nt{Llamada}: las dos alternativas quedan
\nt{Expresion} \tl{;}, la misma producción dos veces. Como conjunto de producciones es una sola, y con
ella hay un solo árbol.

\paso{Paso 2 (\tec{AG}).} \regla{SentenciaExpresion}{\nt{Expresion}}.

\paso{Resultado}
\begin{producciones}
\regla{Sentencia}{\tl{;} \alt{} \nt{SentenciaExpresion} \tl{;} \alt{} \nt{VarLocal} \tl{;} \alt{} \nt{Return} \tl{;} \alt{} \nt{If} \alt{} \nt{While} \alt{} \nt{Bloque}}
\regla{SentenciaExpresion}{\nt{Expresion}}
\end{producciones}
\cj{\tl{;}}, \cE, \cj{\tl{var}}, \cj{\tl{return}}, \cj{\tl{if}}, \cj{\tl{while}} y \cj{\tl{\{}}. \ok{}
Distinguir una asignación de una llamada exige el significado de la expresión: es del semántico.

\subsubsection*{F3 --- \nt{If}}

\paso{Original}
\begin{producciones}
\regla{If}{\tl{if} \tl{(} \nt{Expresion} \tl{)} \nt{Sentencia} \alt{} \tl{if} \tl{(} \nt{Expresion} \tl{)} \nt{Sentencia} \tl{else} \nt{Sentencia}}
\end{producciones}
\paso{Conflicto.} \cj{\tl{if}} contra \cj{\tl{if}}.

\paso{Paso.} \tec{FD} con $\alpha =$ \tl{if} \tl{(} \nt{Expresion} \tl{)} \nt{Sentencia},
$\beta_1 = \epsilon$, $\beta_2 =$ \tl{else} \nt{Sentencia} y $R =$ \nt{ElseOpcional}.

\paso{Resultado}
\begin{producciones}
\regla{If}{\tl{if} \tl{(} \nt{Expresion} \tl{)} \nt{Sentencia} \nt{ElseOpcional}}
\regla{ElseOpcional}{\tl{else} \nt{Sentencia} \alt{} \eps}
\end{producciones}
\cj{\tl{else}} contra \cj{\eps}. \ok{} Queda la ambigüedad (2): \nt{ElseOpcional} deriva a \eps{} y lo
que le sigue puede empezar con \tl{else} (§\ref{sec:ambiguedades}).

\subsubsection*{F4 --- \nt{DimensionesConTamanio}}

\paso{Original}
\begin{producciones}
\regla{DimensionesConTamanio}{\tl{[} \nt{Expresion} \tl{]} \nt{DimensionesConTamanio} \alt{} \tl{[} \nt{Expresion} \tl{]}}
\end{producciones}
\paso{Conflicto.} \cj{\tl{[}} contra \cj{\tl{[}}.

\paso{Paso.} \tec{FD} con $\alpha =$ \tl{[} \nt{Expresion} \tl{]}, $\beta_1 =$
\nt{DimensionesConTamanio}, $\beta_2 = \epsilon$ y $R =$ \nt{DimensionesConTamanioOpcional}.

\paso{Resultado}
\begin{producciones}
\regla{DimensionesConTamanio}{\tl{[} \nt{Expresion} \tl{]} \nt{DimensionesConTamanioOpcional}}
\regla{DimensionesConTamanioOpcional}{\nt{DimensionesConTamanio} \alt{} \eps}
\end{producciones}
\cj{\tl{[}} contra \cj{\eps}. \ok{} Aparece la ambigüedad (3): lo que sigue a
\nt{DimensionesConTamanioOpcional}, por \nt{Primario} y \nt{RestoReferencia}, puede empezar con
\tl{[} (§\ref{sec:ambiguedades}).

\subsubsection*{F5 --- \nt{ListaExps}}

\paso{Original}
\begin{producciones}
\regla{ListaExps}{\nt{Expresion} \alt{} \nt{Expresion} \tl{,} \nt{ListaExps}}
\end{producciones}
\paso{Conflicto.} Las dos alternativas tienen Primeros \cE.

\paso{Paso.} \tec{FD} con $\alpha =$ \nt{Expresion}, $\beta_1 = \epsilon$, $\beta_2 =$ \tl{,}
\nt{ListaExps} y $R =$ \nt{RestoListaExps}.

\paso{Resultado}
\begin{producciones}
\regla{ListaExps}{\nt{Expresion} \nt{RestoListaExps}}
\regla{RestoListaExps}{\tl{,} \nt{ListaExps} \alt{} \eps}
\end{producciones}
\cj{\tl{,}} contra \cj{\eps}. \ok

\subsubsection*{F6 --- \nt{Miembro}}

\paso{Original}
\begin{producciones}
\regla{Miembro}{\nt{Atributo} \alt{} \nt{Metodo} \alt{} \nt{Constructor}}
\regla{Atributo}{\nt{Tipo} \tl{idMetVar} \tl{;}}
\regla{Metodo}{\nt{ModificadorOpcional} \nt{TipoMetodo} \tl{idMetVar} \nt{ArgsFormales} \nt{Bloque}}
\regla{Constructor}{\tl{public} \tl{idClase} \nt{ArgsFormales} \nt{Bloque}}
\regla{ModificadorOpcional}{\tl{static} \alt{} \eps}
\regla{TipoMetodo}{\nt{Tipo} \alt{} \tl{void}}
\end{producciones}
\paso{Conflicto.} $\prim{\nt{Atributo}} = \cT$, $\prim{\nt{Metodo}} = \cj{\tl{static}, \tl{void}}
\cup \cT$ y $\prim{\nt{Constructor}} = \cj{\tl{public}}$: \nt{Atributo} y \nt{Metodo} chocan en \cT.
El prefijo no es directo.

\paso{Paso 1.} \tec{SU} de \nt{Atributo} y \nt{Metodo}:
\begin{producciones}
\regla{Miembro}{\nt{Tipo} \tl{idMetVar} \tl{;} \alt{} \nt{ModificadorOpcional} \nt{TipoMetodo} \tl{idMetVar} \nt{ArgsFormales} \nt{Bloque} \alt{} \nt{Constructor}}
\end{producciones}
\paso{Paso 2.} \tec{SU} de \nt{ModificadorOpcional}:
\begin{producciones}
\regla{Miembro}{\nt{Tipo} \tl{idMetVar} \tl{;}}
\regla{Miembro}{\tl{static} \nt{TipoMetodo} \tl{idMetVar} \nt{ArgsFormales} \nt{Bloque}}
\regla{Miembro}{\nt{TipoMetodo} \tl{idMetVar} \nt{ArgsFormales} \nt{Bloque}}
\regla{Miembro}{\nt{Constructor}}
\end{producciones}
\paso{Paso 3.} \tec{SU} de \nt{TipoMetodo} en la tercera:
\begin{producciones}
\regla{Miembro}{\nt{Tipo} \tl{idMetVar} \tl{;}}
\regla{Miembro}{\tl{static} \nt{TipoMetodo} \tl{idMetVar} \nt{ArgsFormales} \nt{Bloque}}
\regla{Miembro}{\nt{Tipo} \tl{idMetVar} \nt{ArgsFormales} \nt{Bloque}}
\regla{Miembro}{\tl{void} \tl{idMetVar} \nt{ArgsFormales} \nt{Bloque}}
\regla{Miembro}{\nt{Constructor}}
\end{producciones}
\paso{Paso 4.} \tec{FD} con $\alpha =$ \nt{Tipo} \tl{idMetVar}, $\beta_1 =$ \tl{;}, $\beta_2 =$
\nt{ArgsFormales} \nt{Bloque} y $R =$ \nt{AtributoOMetodo}:
\begin{producciones}
\regla{Miembro}{\nt{Tipo} \tl{idMetVar} \nt{AtributoOMetodo}}
\regla{Miembro}{\tl{static} \nt{TipoMetodo} \tl{idMetVar} \nt{ArgsFormales} \nt{Bloque}}
\regla{Miembro}{\tl{void} \tl{idMetVar} \nt{ArgsFormales} \nt{Bloque}}
\regla{Miembro}{\nt{Constructor}}
\regla{AtributoOMetodo}{\tl{;} \alt{} \nt{ArgsFormales} \nt{Bloque}}
\end{producciones}
\nt{Miembro}: \cT, \cj{\tl{static}}, \cj{\tl{void}} y \cj{\tl{public}}. \nt{AtributoOMetodo}:
\cj{\tl{;}} contra \cj{\tl{(}}. \ok

\paso{Paso 5 (\tec{AG}).} \regla{CuerpoDeMetodo}{\nt{ArgsFormales} \nt{Bloque}} en sus cuatro
apariciones, incluida la de \nt{Constructor}, y \regla{Metodo}{\tl{idMetVar} \nt{CuerpoDeMetodo}} en
las ramas \tl{static} y \tl{void}.

\paso{Resultado}
\begin{producciones}
\regla{Miembro}{\nt{Tipo} \tl{idMetVar} \nt{AtributoOMetodo} \alt{} \tl{static} \nt{TipoMetodo} \nt{Metodo} \alt{} \tl{void} \nt{Metodo} \alt{} \nt{Constructor}}
\regla{AtributoOMetodo}{\tl{;} \alt{} \nt{CuerpoDeMetodo}}
\regla{Metodo}{\tl{idMetVar} \nt{CuerpoDeMetodo}}
\regla{CuerpoDeMetodo}{\nt{ArgsFormales} \nt{Bloque}}
\regla{Constructor}{\tl{public} \tl{idClase} \nt{CuerpoDeMetodo}}
\end{producciones}
Se pierden \nt{Atributo} y \nt{ModificadorOpcional}. \nt{Metodo} conserva el nombre pero deriva la
cola de las ramas \tl{static} y \tl{void}; el método con tipo y sin \tl{static} llega por
\nt{AtributoOMetodo}, porque su \tl{idMetVar} se consumió al factorizar contra el atributo.

\subsubsection*{F7 --- \nt{Primario}, prefijo \tl{idMetVar}}

\paso{Original}
\begin{producciones}
\regla{Primario}{\tl{this} \alt{} \tl{stringLiteral} \alt{} \nt{AccesoVar} \alt{} \nt{CrearArreglo} \alt{} \nt{LlamadaConstructor} \alt{} \nt{LlamadaMetodo} \alt{} \nt{LlamadaMetodoEstatico} \alt{} \nt{ExpresionParentizada}}
\regla{AccesoVar}{\tl{idMetVar}}
\regla{LlamadaMetodo}{\tl{idMetVar} \nt{ArgsActuales}}
\end{producciones}
\paso{Conflicto.} \nt{AccesoVar} y \nt{LlamadaMetodo} en \cj{\tl{idMetVar}}. El conflicto en
\cj{\tl{new}} se trata en F8.

\paso{Paso 1 (\tec{AG}).} \regla{AccesoVarOLlamadaMetodo}{\nt{AccesoVar} \alt{} \nt{LlamadaMetodo}}
reemplaza a las dos en \nt{Primario}.

\paso{Paso 2.} \tec{FP}: \tec{SU} de las dos, que da \regla{AccesoVarOLlamadaMetodo}{\tl{idMetVar}
\alt{} \tl{idMetVar} \nt{ArgsActuales}}, y \tec{FD} con $\alpha =$ \tl{idMetVar}, $\beta_1 =
\epsilon$, $\beta_2 =$ \nt{ArgsActuales} y $R =$ \nt{AccesoOLlamada}.

\paso{Resultado}
\begin{producciones}
\regla{AccesoVarOLlamadaMetodo}{\tl{idMetVar} \nt{AccesoOLlamada}}
\regla{AccesoOLlamada}{\nt{ArgsActuales} \alt{} \eps}
\end{producciones}
\cj{\tl{(}} contra \cj{\eps}. \ok{} No se comparte con \nt{EncadenadoAtributoOMetodo} aunque tengan la
misma forma: su \eps{} es una variable propia, y el de aquél un atributo de otro objeto.

\subsubsection*{F8 --- \nt{Primario}, prefijo \tl{new}}

\paso{Original}
\begin{producciones}
\regla{CrearArreglo}{\tl{new} \nt{TipoBase} \nt{DimensionesConTamanio}}
\regla{LlamadaConstructor}{\tl{new} \nt{TipoReferencia} \nt{ArgsActuales}}
\regla{TipoBase}{\nt{TipoPrimitivo} \alt{} \nt{TipoReferencia} \alt{} \tl{idGen}}
\end{producciones}
\paso{Conflicto.} \cj{\tl{new}} contra \cj{\tl{new}}.

\paso{Paso 1 (\tec{AG}).} \regla{Creacion}{\nt{CrearArreglo} \alt{} \nt{LlamadaConstructor}} reemplaza
a las dos en \nt{Primario}.

\paso{Paso 2.} \tec{FP}: \tec{SU} de las dos y \tec{FD} con $\alpha =$ \tl{new} y $R =$
\nt{ArregloOInstancia}:
\begin{producciones}
\regla{Creacion}{\tl{new} \nt{ArregloOInstancia}}
\regla{ArregloOInstancia}{\nt{TipoBase} \nt{DimensionesConTamanio} \alt{} \nt{TipoReferencia} \nt{ArgsActuales}}
\end{producciones}
Conflicto nuevo: \cT{} contra \cj{\tl{idClase}}.

\paso{Paso 3.} \tec{FP}: \tec{SU} de \nt{TipoBase},
\begin{producciones}
\regla{ArregloOInstancia}{\nt{TipoPrimitivo} \nt{DimensionesConTamanio} \alt{} \nt{TipoReferencia} \nt{DimensionesConTamanio} \alt{} \tl{idGen} \nt{DimensionesConTamanio} \alt{} \nt{TipoReferencia} \nt{ArgsActuales}}
\end{producciones}
y \tec{FD} con $\alpha =$ \nt{TipoReferencia}, $\beta_1 =$ \nt{DimensionesConTamanio}, $\beta_2 =$
\nt{ArgsActuales} y $R =$ \nt{ArregloOInstanciaDeClase}. \nt{TipoReferencia} no se sube: no hace
falta.

\paso{Resultado}
\begin{producciones}
\regla{Primario}{\tl{this} \alt{} \tl{stringLiteral} \alt{} \nt{ExpresionParentizada} \alt{} \nt{AccesoVarOLlamadaMetodo} \alt{} \nt{LlamadaMetodoEstatico} \alt{} \nt{Creacion}}
\regla{Creacion}{\tl{new} \nt{ArregloOInstancia}}
\regla{ArregloOInstancia}{\nt{TipoPrimitivo} \nt{DimensionesConTamanio} \alt{} \tl{idGen} \nt{DimensionesConTamanio} \alt{} \nt{TipoReferencia} \nt{ArregloOInstanciaDeClase}}
\regla{ArregloOInstanciaDeClase}{\nt{DimensionesConTamanio} \alt{} \nt{ArgsActuales}}
\end{producciones}
\nt{Primario}: \cj{\tl{this}}, \cj{\tl{stringLiteral}}, \cj{\tl{(}}, \cj{\tl{idMetVar}},
\cj{\tl{idClase}} y \cj{\tl{new}}. \nt{ArregloOInstancia}: \cj{\tl{boolean}, \tl{char}, \tl{int}},
\cj{\tl{idGen}} y \cj{\tl{idClase}}. \nt{ArregloOInstanciaDeClase}: \cj{\tl{[}} contra \cj{\tl{(}}.
\ok{} Se pierden \nt{AccesoVar}, \nt{LlamadaMetodo}, \nt{CrearArreglo} y \nt{LlamadaConstructor}.

\subsection{El núcleo}

Con R1--R3 y F1--F8 la gramática es LL(1): de 60 a 67 no terminales y de 126 a 135 producciones.
Genera el mismo lenguaje que la de la cátedra. Los diez no terminales que se pierden lo hacen por
\tec{SU}.

% ---------------------------------------------------------------------------
\subsection{Logros}

Cada logro parte de la gramática que dejó el anterior: primero se escribe su extensión y después se
la lleva a LL(1) con las mismas técnicas.

\subsubsection*{L1 --- Floats (Etapa 1)}
\label{sec:l1}

\paso{Extensión}
\begin{producciones}
\regla{TipoPrimitivo}{\tl{boolean} \alt{} \tl{char} \alt{} \tl{int} \alt{} \tl{float}}
\regla{Primitivo}{\tl{true} \alt{} \tl{false} \alt{} \tl{intLiteral} \alt{} \tl{charLiteral} \alt{} \tl{null} \alt{} \tl{floatLiteral}}
\end{producciones}
Sin conflictos.

\subsubsection*{L2 --- Operadores Posifijos}
\label{sec:l2}

\paso{Extensión}
\begin{producciones}
\regla{Operando}{\nt{Primitivo} \alt{} \nt{Referencia} \nt{PosfijoOpcional}}
\regla{PosfijoOpcional}{\tl{++} \alt{} \tl{-{}-} \alt{} \eps}
\end{producciones}
\nt{Operando}: los literales contra \cP. \nt{PosfijoOpcional}: \cj{\tl{++}}, \cj{\tl{-{}-}} y
\cj{\eps}. Sin conflictos.

\subsubsection*{L3 --- Operador Ternario}
\label{sec:l3}

\paso{Extensión} Sobre las reglas de la base, el lado derecho del \tl{=} y el de \nt{VarLocal} pasan
a ser una \nt{ExpresionCondicional}, con los operandos del JLS §15.25:
\begin{producciones}
\regla{Expresion}{\nt{ExpresionCompuesta} \nt{OperadorAsignacion} \nt{ExpresionCondicional} \alt{} \nt{ExpresionCondicional}}
\regla{ExpresionCondicional}{\nt{ExpresionCompuesta} \tl{?} \nt{Expresion} \tl{:} \nt{ExpresionCondicional} \alt{} \nt{ExpresionCompuesta}}
\regla{VarLocal}{\tl{var} \tl{idMetVar} \tl{=} \nt{ExpresionCondicional}}
\end{producciones}
\paso{Conflictos.} En \nt{ExpresionCondicional} y en \nt{Expresion} las dos alternativas tienen
Primeros \cE.

\paso{Paso 1.} \tec{FD} en \nt{ExpresionCondicional} con $\alpha =$ \nt{ExpresionCompuesta} y $R =$
\nt{RestoCondicional}.

\paso{Paso 2.} \tec{FP} en \nt{Expresion}: \tec{SU} de \nt{ExpresionCondicional} en la segunda, que
queda \nt{ExpresionCompuesta} \nt{RestoCondicional}, y \tec{FD} con $\alpha =$
\nt{ExpresionCompuesta}, $\beta_1 =$ \nt{OperadorAsignacion} \nt{ExpresionCondicional}, $\beta_2 =$
\nt{RestoCondicional} y $R =$ \nt{RestoExpresion}, que reemplaza al de F1.

\paso{Resultado}
\begin{producciones}
\regla{Expresion}{\nt{ExpresionCompuesta} \nt{RestoExpresion}}
\regla{RestoExpresion}{\nt{OperadorAsignacion} \nt{ExpresionCondicional} \alt{} \nt{RestoCondicional}}
\regla{ExpresionCondicional}{\nt{ExpresionCompuesta} \nt{RestoCondicional}}
\regla{RestoCondicional}{\tl{?} \nt{Expresion} \tl{:} \nt{ExpresionCondicional} \alt{} \eps}
\end{producciones}
\nt{RestoExpresion}: \cj{\tl{=}} contra \cj{\tl{?}, \eps}. \nt{RestoCondicional}: \cj{\tl{?}} contra
\cj{\eps}. \ok{} El lado izquierdo del \tl{=} sigue siendo una \nt{ExpresionCompuesta} y el derecho no
admite otro \tl{=}. El ternario es asociativo a derecha.

\subsubsection*{L4 --- Variables Locales Clásicas}
\label{sec:l4}

\paso{Extensión}
\begin{producciones}
\regla{Sentencia}{\dots{} \alt{} \nt{VarLocalClasica} \tl{;}}
\regla{VarLocalClasica}{\nt{Tipo} \nt{ListaDeclaradores}}
\regla{ListaDeclaradores}{\tl{idMetVar} \nt{InicializacionDeDeclarador} \alt{} \nt{ListaDeclaradores} \tl{,} \tl{idMetVar} \nt{InicializacionDeDeclarador}}
\regla{InicializacionDeDeclarador}{\tl{=} \nt{ExpresionCondicional} \alt{} \eps}
\end{producciones}
\paso{Conflictos.} \nt{ListaDeclaradores} tiene recursión a izquierda. En \nt{Sentencia},
\nt{SentenciaExpresion} \tl{;} (\cE) y \nt{VarLocalClasica} \tl{;} (\cT) chocan en
\cj{\tl{idClase}}: \texttt{C.m();} contra \texttt{C x;}.

\paso{Paso 1.} \tec{RI} en \nt{ListaDeclaradores} con $\beta_1 =$ \tl{idMetVar}
\nt{InicializacionDeDeclarador}, $\alpha_1 =$ \tl{,} \tl{idMetVar} \nt{InicializacionDeDeclarador} y
$R =$ \nt{RestoListaDeclaradores}.

\paso{Paso 2.} \tec{FP} en \nt{Sentencia}: \tec{SU} hasta exponer \tl{idClase}. La declaración se
sube por \nt{VarLocalClasica}, \nt{Tipo}, \nt{TipoBase} y \nt{TipoReferencia}; la expresión, por
\nt{SentenciaExpresion}, \nt{Expresion}, \nt{ExpresionCompuesta}, \nt{ExpresionBasica},
\nt{Operando}, \nt{Referencia}, \nt{Primario} y \nt{LlamadaMetodoEstatico}. Con $X$ cada uno de
\tl{this}, \tl{stringLiteral}, \nt{ExpresionParentizada}, \nt{AccesoVarOLlamadaMetodo} y
\nt{Creacion}, esas dos alternativas pasan a ser:
\begin{producciones}
\regla{Sentencia}{\nt{TipoPrimitivo} \nt{DimensionesOpcionales} \nt{ListaDeclaradores} \tl{;}}
\regla{Sentencia}{\tl{idGen} \nt{DimensionesOpcionales} \nt{ListaDeclaradores} \tl{;}}
\regla{Sentencia}{\tl{idClase} \nt{TipoGenericoOpcional} \nt{DimensionesOpcionales} \nt{ListaDeclaradores} \tl{;}}
\regla{Sentencia}{\nt{OperadorUnario} \nt{Operando} \nt{RestoExpresionCompuesta} \nt{RestoExpresion} \tl{;}}
\regla{Sentencia}{\nt{Primitivo} \nt{RestoExpresionCompuesta} \nt{RestoExpresion} \tl{;}}
\regla{Sentencia}{$X$ \nt{RestoReferencia} \nt{PosfijoOpcional} \nt{RestoExpresionCompuesta} \nt{RestoExpresion} \tl{;}}
\regla{Sentencia}{\tl{idClase} \tl{.} \tl{idMetVar} \nt{ArgsActuales} \nt{RestoReferencia} \nt{PosfijoOpcional} \nt{RestoExpresionCompuesta} \nt{RestoExpresion} \tl{;}}
\end{producciones}
\paso{Paso 3.} \tec{FD} con $\alpha =$ \tl{idClase} y $R =$ \nt{SentenciaDesdeIdClase}.

\paso{Paso 4 (\tec{AG}).} Para no duplicar las colas:
\begin{producciones}
\regla{ColaTrasOperando}{\nt{RestoExpresionCompuesta} \nt{RestoExpresion}}
\regla{ColaDeExpresionDesdePrimario}{\nt{RestoReferencia} \nt{PosfijoOpcional} \nt{ColaTrasOperando}}
\regla{PrimarioSinIdClase}{\tl{this} \alt{} \tl{stringLiteral} \alt{} \nt{ExpresionParentizada} \alt{} \nt{AccesoVarOLlamadaMetodo} \alt{} \nt{Creacion}}
\regla{Primario}{\nt{LlamadaMetodoEstatico} \alt{} \nt{PrimarioSinIdClase}}
\regla{ExpresionSinIdClase}{\nt{OperadorUnario} \nt{Operando} \nt{ColaTrasOperando} \alt{} \nt{Primitivo} \nt{ColaTrasOperando} \alt{} \nt{PrimarioSinIdClase} \nt{ColaDeExpresionDesdePrimario}}
\end{producciones}
\paso{Resultado}
\begin{producciones}
\regla{Sentencia}{\tl{;} \alt{} \nt{VarLocal} \tl{;} \alt{} \nt{Return} \tl{;} \alt{} \nt{If} \alt{} \nt{While} \alt{} \nt{Bloque} \alt{} \nt{ExpresionSinIdClase} \tl{;}}
\regla{Sentencia}{\nt{TipoPrimitivo} \nt{DimensionesOpcionales} \nt{ListaDeclaradores} \tl{;} \alt{} \tl{idGen} \nt{DimensionesOpcionales} \nt{ListaDeclaradores} \tl{;} \alt{} \tl{idClase} \nt{SentenciaDesdeIdClase}}
\regla{SentenciaDesdeIdClase}{\nt{TipoGenericoOpcional} \nt{DimensionesOpcionales} \nt{ListaDeclaradores} \tl{;} \alt{} \tl{.} \tl{idMetVar} \nt{ArgsActuales} \nt{ColaDeExpresionDesdePrimario} \tl{;}}
\regla{ListaDeclaradores}{\tl{idMetVar} \nt{InicializacionDeDeclarador} \nt{RestoListaDeclaradores}}
\regla{RestoListaDeclaradores}{\tl{,} \tl{idMetVar} \nt{InicializacionDeDeclarador} \nt{RestoListaDeclaradores} \alt{} \eps}
\end{producciones}
más los cinco no terminales del paso 4. \nt{Sentencia}: \cj{\tl{;}}, \cj{\tl{var}}, \cj{\tl{return}},
\cj{\tl{if}}, \cj{\tl{while}}, \cj{\tl{\{}}, \cE{} sin \tl{idClase}, \cj{\tl{boolean}, \tl{char},
\tl{float}, \tl{int}}, \cj{\tl{idGen}} y \cj{\tl{idClase}}. \nt{SentenciaDesdeIdClase}: \cj{\tl{<},
\tl{[}, \tl{idMetVar}} contra \cj{\tl{.}}. \nt{Primario}: \cj{\tl{idClase}} contra el resto de \cP.
\ok{} Se disuelven \nt{SentenciaExpresion} y \nt{VarLocalClasica}. Mirar dos tokens después del
\tl{idClase} evitaría los no terminales nuevos, pero se aparta del Esquema Simple, que usa uno.

\subsubsection*{L5 --- Atributos Inicializados}
\label{sec:l5}

\paso{Extensión} La cola del atributo admite la inicialización:
\begin{producciones}
\regla{AtributoOMetodo}{\nt{InicializacionDeAtributo} \tl{;} \alt{} \nt{CuerpoDeMetodo}}
\regla{InicializacionDeAtributo}{\tl{=} \nt{ExpresionCondicional} \alt{} \eps}
\end{producciones}
\cj{\tl{=}, \tl{;}} contra \cj{\tl{(}}. Sin conflictos. No se comparte con
\nt{InicializacionDeDeclarador}: para el semántico son construcciones distintas.

\subsubsection*{L6 --- Visibilidad Mejorada}
\label{sec:l6}

\paso{Extensión} Visibilidad opcional en atributos, métodos y constructores; \tl{static} antes o
después de ella, sin repetir ninguno; \tl{public} opcional en los métodos de interfaz:
\begin{producciones}
\regla{Miembro}{\nt{VisibilidadOpcional} \nt{Tipo} \tl{idMetVar} \nt{AtributoOMetodo} \alt{} \nt{Visibilidad} \tl{static} \nt{TipoMetodo} \nt{Metodo}}
\regla{Miembro}{\tl{static} \nt{VisibilidadOpcional} \nt{TipoMetodo} \nt{Metodo} \alt{} \nt{VisibilidadOpcional} \tl{void} \nt{Metodo} \alt{} \nt{Constructor}}
\regla{Constructor}{\nt{VisibilidadOpcional} \tl{idClase} \nt{CuerpoDeMetodo}}
\regla{Visibilidad}{\tl{public} \alt{} \tl{private}}
\regla{VisibilidadOpcional}{\nt{Visibilidad} \alt{} \eps}
\regla{MetodoInterfaz}{\nt{PublicOpcional} \nt{TipoMetodo} \tl{idMetVar} \nt{ArgsFormales} \tl{;}}
\regla{PublicOpcional}{\tl{public} \alt{} \eps}
\end{producciones}
\paso{Conflictos.} En \nt{Miembro}, \cj{\tl{public}, \tl{private}} en la 1ª, 2ª, 4ª y 5ª, e
\cj{\tl{idClase}} en la 1ª y la 5ª.

\paso{Paso 1.} \tec{SU} de \nt{VisibilidadOpcional} en la 1ª y la 4ª, y de \nt{Constructor}:
\begin{producciones}
\regla{Miembro}{\nt{Visibilidad} \nt{Tipo} \tl{idMetVar} \nt{AtributoOMetodo} \alt{} \nt{Tipo} \tl{idMetVar} \nt{AtributoOMetodo}}
\regla{Miembro}{\nt{Visibilidad} \tl{static} \nt{TipoMetodo} \nt{Metodo} \alt{} \tl{static} \nt{VisibilidadOpcional} \nt{TipoMetodo} \nt{Metodo}}
\regla{Miembro}{\nt{Visibilidad} \tl{void} \nt{Metodo} \alt{} \tl{void} \nt{Metodo}}
\regla{Miembro}{\nt{Visibilidad} \tl{idClase} \nt{CuerpoDeMetodo} \alt{} \tl{idClase} \nt{CuerpoDeMetodo}}
\end{producciones}
\paso{Paso 2 (\tec{AG}).} Las tres construcciones que aparecen con y sin visibilidad:
\begin{producciones}
\regla{MiembroSinModificadores}{\tl{void} \nt{Metodo} \alt{} \nt{Tipo} \tl{idMetVar} \nt{AtributoOMetodo} \alt{} \tl{idClase} \nt{CuerpoDeMetodo}}
\regla{Miembro}{\nt{Visibilidad} \nt{MiembroSinModificadores} \alt{} \nt{Visibilidad} \tl{static} \nt{TipoMetodo} \nt{Metodo} \alt{} \tl{static} \nt{VisibilidadOpcional} \nt{TipoMetodo} \nt{Metodo} \alt{} \nt{MiembroSinModificadores}}
\end{producciones}
\paso{Paso 3.} \tec{FD} en \nt{Miembro} con $\alpha =$ \nt{Visibilidad} y $R =$
\nt{MiembroTrasVisibilidad}.

\paso{Paso 4.} En \nt{MiembroSinModificadores} chocan \cT{} y \cj{\tl{idClase}}: \tec{FP} con
\tec{SU} de \nt{Tipo}, \nt{TipoBase} y \nt{TipoReferencia}, y \tec{FD} con $\alpha =$ \tl{idClase} y
$R =$ \nt{MiembroDesdeIdClase}.

\paso{Resultado}
\begin{producciones}
\regla{Miembro}{\nt{Visibilidad} \nt{MiembroTrasVisibilidad} \alt{} \tl{static} \nt{VisibilidadOpcional} \nt{TipoMetodo} \nt{Metodo} \alt{} \nt{MiembroSinModificadores}}
\regla{MiembroTrasVisibilidad}{\tl{static} \nt{TipoMetodo} \nt{Metodo} \alt{} \nt{MiembroSinModificadores}}
\regla{MiembroSinModificadores}{\tl{void} \nt{Metodo} \alt{} \nt{TipoPrimitivo} \nt{DimensionesOpcionales} \tl{idMetVar} \nt{AtributoOMetodo}}
\regla{MiembroSinModificadores}{\tl{idGen} \nt{DimensionesOpcionales} \tl{idMetVar} \nt{AtributoOMetodo} \alt{} \tl{idClase} \nt{MiembroDesdeIdClase}}
\regla{MiembroDesdeIdClase}{\nt{CuerpoDeMetodo} \alt{} \nt{TipoGenericoOpcional} \nt{DimensionesOpcionales} \tl{idMetVar} \nt{AtributoOMetodo}}
\end{producciones}
\nt{Miembro}: \cj{\tl{public}, \tl{private}}, \cj{\tl{static}} y \cj{\tl{boolean}, \tl{char},
\tl{float}, \tl{int}, \tl{idClase}, \tl{idGen}, \tl{void}}. \nt{MiembroTrasVisibilidad}:
\cj{\tl{static}} contra ese último. \nt{MiembroSinModificadores}: \cj{\tl{void}}, \cj{\tl{boolean},
\tl{char}, \tl{float}, \tl{int}}, \cj{\tl{idGen}} y \cj{\tl{idClase}}. \nt{MiembroDesdeIdClase}:
\cj{\tl{(}} contra \cj{\tl{<}, \tl{[}, \tl{idMetVar}}. \ok{} Se disuelve \nt{Constructor}: sin un
\tl{public} obligatorio, compite con el tipo por \tl{idClase}.

\subsubsection*{L7 --- Fors!}
\label{sec:l7}

\paso{Extensión}
\begin{producciones}
\regla{Sentencia}{\dots{} \alt{} \nt{For}}
\regla{For}{\tl{for} \tl{(} \nt{EncabezadoFor} \tl{)} \nt{Sentencia}}
\regla{EncabezadoFor}{\nt{VarLocal} \nt{ClausulasFinalesFor} \alt{} \nt{Expresion} \nt{ClausulasFinalesFor} \alt{} \tl{var} \tl{idMetVar} \tl{:} \nt{Expresion}}
\regla{ClausulasFinalesFor}{\tl{;} \nt{Expresion} \tl{;} \nt{Expresion}}
\end{producciones}
\paso{Conflicto.} En \nt{EncabezadoFor}, \cj{\tl{var}} contra \cj{\tl{var}}.

\paso{Paso.} \tec{FP}: \tec{SU} de \nt{VarLocal}, que da \tl{var} \tl{idMetVar} \tl{=}
\nt{ExpresionCondicional} \nt{ClausulasFinalesFor}, y \tec{FD} con $\alpha =$ \tl{var} \tl{idMetVar} y
$R =$ \nt{EncabezadoForTrasVar}.

\paso{Resultado}
\begin{producciones}
\regla{EncabezadoFor}{\tl{var} \tl{idMetVar} \nt{EncabezadoForTrasVar} \alt{} \nt{Expresion} \nt{ClausulasFinalesFor}}
\regla{EncabezadoForTrasVar}{\tl{:} \nt{Expresion} \alt{} \tl{=} \nt{ExpresionCondicional} \nt{ClausulasFinalesFor}}
\end{producciones}
\nt{EncabezadoFor}: \cj{\tl{var}} contra \cE. \nt{EncabezadoForTrasVar}: \cj{\tl{:}} contra
\cj{\tl{=}}. \ok

\subsubsection*{L8 --- Inicialización de Arreglos}
\label{sec:l8}

\paso{Extensión} Donde F8 dejó \nt{DimensionesConTamanio} va una \nt{CreacionDeArreglo}, que admite
además la forma de Java con inicializador (JLS §15.10.1 y §10.6):
\begin{producciones}
\regla{ArregloOInstancia}{\nt{TipoPrimitivo} \nt{CreacionDeArreglo} \alt{} \tl{idGen} \nt{CreacionDeArreglo} \alt{} \nt{TipoReferencia} \nt{ArregloOInstanciaDeClase}}
\regla{ArregloOInstanciaDeClase}{\nt{CreacionDeArreglo} \alt{} \nt{ArgsActuales}}
\regla{CreacionDeArreglo}{\nt{DimensionesConTamanio} \alt{} \tl{[} \tl{]} \nt{DimensionesOpcionales} \nt{InicializadorDeArreglo}}
\regla{InicializadorDeArreglo}{\tl{\{} \nt{ListaElementosOpcional} \tl{\}}}
\regla{ListaElementosOpcional}{\nt{ListaElementos} \alt{} \tl{,} \alt{} \eps}
\regla{ListaElementos}{\nt{ElementoInicializador} \alt{} \nt{ElementoInicializador} \tl{,} \alt{} \nt{ElementoInicializador} \tl{,} \nt{ListaElementos}}
\regla{ElementoInicializador}{\nt{Expresion} \alt{} \nt{InicializadorDeArreglo}}
\end{producciones}
\paso{Conflictos.} En \nt{CreacionDeArreglo}, \cj{\tl{[}} contra \cj{\tl{[}}. En \nt{ListaElementos},
las tres alternativas tienen Primeros $\cE \cup \cj{\tl{\{}}$.

\paso{Paso 1.} \tec{FP} en \nt{CreacionDeArreglo}: \tec{SU} de \nt{DimensionesConTamanio}, que da
\tl{[} \nt{Expresion} \tl{]} \nt{DimensionesConTamanioOpcional}, y \tec{FD} con $\alpha =$ \tl{[} y
$R =$ \nt{RestoCreacionDeArreglo}.

\paso{Paso 2.} \tec{FD} en \nt{ListaElementos} con $\alpha =$ \nt{ElementoInicializador} y $R =$
\nt{RestoListaElementos}, que queda \regla{RestoListaElementos}{\eps{} \alt{} \tl{,} \alt{} \tl{,}
\nt{ListaElementos}}: conflicto \cj{\tl{,}} contra \cj{\tl{,}}.

\paso{Paso 3.} \tec{FD} en \nt{RestoListaElementos} con $\alpha =$ \tl{,}, $\beta_1 = \epsilon$,
$\beta_2 =$ \nt{ListaElementos} y $R =$ \nt{ElementosTrasComa}.

\paso{Resultado}
\begin{producciones}
\regla{CreacionDeArreglo}{\tl{[} \nt{RestoCreacionDeArreglo}}
\regla{RestoCreacionDeArreglo}{\nt{Expresion} \tl{]} \nt{DimensionesConTamanioOpcional} \alt{} \tl{]} \nt{DimensionesOpcionales} \nt{InicializadorDeArreglo}}
\regla{ListaElementos}{\nt{ElementoInicializador} \nt{RestoListaElementos}}
\regla{RestoListaElementos}{\tl{,} \nt{ElementosTrasComa} \alt{} \eps}
\regla{ElementosTrasComa}{\nt{ListaElementos} \alt{} \eps}
\end{producciones}
\nt{RestoCreacionDeArreglo}: \cE{} contra \cj{\tl{]}}. \nt{ListaElementosOpcional}: $\cE \cup
\cj{\tl{\{}}$, \cj{\tl{,}} y \cj{\eps}. \nt{RestoListaElementos}: \cj{\tl{,}} contra \cj{\eps}.
\nt{ElementosTrasComa}: $\cE \cup \cj{\tl{\{}}$ contra \cj{\eps}. \nt{ElementoInicializador}: \cE{}
contra \cj{\tl{\{}}. \ok{} Sin inicializador sigue la creación de la base, con toda dimensión con su
tamaño. \nt{InicializadorDeArreglo} sólo es alcanzable desde un \tl{new}.

\subsubsection*{L9 --- Genericidad Avanzada}
\label{sec:l9}

\paso{Extensión, anidamiento.} \regla{InstanciadoOParametrico}{\tl{idGen} \alt{} \tl{idClase}
\nt{TipoGenericoOpcional}}. \cj{\tl{idGen}} contra \cj{\tl{idClase}}: sin conflictos.

\paso{Extensión, diamante.} Sólo al instanciar: \regla{ArregloOInstancia}{\dots{} \alt{} \tl{idClase}
\tl{<} \tl{>} \nt{ArgsActuales}}.

\paso{Conflicto.} \nt{TipoReferencia} \nt{ArregloOInstanciaDeClase} contra \tl{idClase} \tl{<}
\tl{>} \nt{ArgsActuales}, en \cj{\tl{idClase}}.

\paso{Paso 1.} \tec{FP}: \tec{SU} de \nt{TipoReferencia}, que da \tl{idClase}
\nt{TipoGenericoOpcional} \nt{ArregloOInstanciaDeClase}, y \tec{FD} con $\alpha =$ \tl{idClase} y $R =$
\nt{CreacionDesdeIdClase}, que queda \regla{CreacionDesdeIdClase}{\nt{TipoGenericoOpcional}
\nt{ArregloOInstanciaDeClase} \alt{} \tl{<} \tl{>} \nt{ArgsActuales}}: conflicto en \cj{\tl{<}}.

\paso{Paso 2.} \tec{FP}: \tec{SU} de \nt{TipoGenericoOpcional}, que da \tl{<}
\nt{InstanciadoOParametrico} \tl{>} \nt{ArregloOInstanciaDeClase} \alt{} \nt{ArregloOInstanciaDeClase},
y \tec{FD} con $\alpha =$ \tl{<} y $R =$ \nt{GenericoDeCreacion}.

\paso{Resultado}
\begin{producciones}
\regla{ArregloOInstancia}{\nt{TipoPrimitivo} \nt{CreacionDeArreglo} \alt{} \tl{idGen} \nt{CreacionDeArreglo} \alt{} \tl{idClase} \nt{CreacionDesdeIdClase}}
\regla{CreacionDesdeIdClase}{\tl{<} \nt{GenericoDeCreacion} \alt{} \nt{ArregloOInstanciaDeClase}}
\regla{GenericoDeCreacion}{\tl{>} \nt{ArgsActuales} \alt{} \nt{InstanciadoOParametrico} \tl{>} \nt{ArregloOInstanciaDeClase}}
\end{producciones}
\nt{CreacionDesdeIdClase}: \cj{\tl{<}} contra \cj{\tl{[}, \tl{(}}. \nt{GenericoDeCreacion}:
\cj{\tl{>}} contra \cj{\tl{idClase}, \tl{idGen}}. \ok{} \texttt{List<> x;} y \texttt{new C<>[3]} no
derivan: el diamante sólo lleva a \nt{ArgsActuales}.

% ---------------------------------------------------------------------------
\subsection{Ambigüedades que quedan}
\label{sec:ambiguedades}

La (1) se eliminó en F2. Las otras dos se resuelven por regla de implementación, y de la misma
forma: el analizador toma la alternativa no vacía apenas ve su token.

\begin{center}\small
\begin{tabular}{@{}>{\raggedright\arraybackslash}p{0.40\linewidth}>{\raggedright\arraybackslash}p{0.25\linewidth}>{\raggedright\arraybackslash}p{0.29\linewidth}@{}}
\toprule
\textbf{No terminal} & \textbf{Caso} & \textbf{Resolución} \\
\midrule
\nt{ElseOpcional} & \texttt{if (a) if (b) s1; else s2;} & el \tl{else} liga al \tl{if} más cercano \\
\nt{DimensionesConTamanioOpcional} & \texttt{new int[3][4]} & el \tl{[} es otra dimensión \\
\bottomrule
\end{tabular}
\end{center}

Las dos coinciden con Java. Detectarlas exige mirar qué sigue a un no terminal anulable; ese cálculo
es una auditoría hecha fuera del compilador, que decide sólo con Primeros.

% ===========================================================================
\section{Gramática final}
\label{sec:gramatica-final}
% ===========================================================================

\begin{producciones}
\regla{Inicial}{\nt{ListaClases} \tl{eof}}
\regla{ListaClases}{\nt{Clase} \nt{ListaClases} \alt{} \nt{Interfaz} \nt{ListaClases} \alt{} \eps}
\regla{Clase}{\tl{class} \tl{idClase} \nt{GenericidadOpcional} \nt{HerenciaOpcional} \tl{\{} \nt{ListaMiembros} \tl{\}}}
\regla{Interfaz}{\tl{interface} \tl{idClase} \nt{GenericidadOpcional} \nt{ExtensionOpcional} \tl{\{} \nt{ListaMetodosInterfaz} \tl{\}}}
\end{producciones}

\noindent{\small\color{grisMedio}Una Clase o Interfaz sigue limitada a UN único parámetro de tipo, también con el logro de genericidad avanzada.\par}

\begin{producciones}
\regla{GenericidadOpcional}{\tl{<} \tl{idGen} \tl{>} \alt{} \eps}
\regla{HerenciaOpcional}{\tl{extends} \nt{TipoReferencia} \alt{} \tl{implements} \nt{TipoReferencia} \alt{} \eps}
\regla{ExtensionOpcional}{\tl{extends} \nt{TipoReferencia} \alt{} \eps}
\regla{ListaMiembros}{\nt{Miembro} \nt{ListaMiembros} \alt{} \eps}
\regla{ListaMetodosInterfaz}{\nt{MetodoInterfaz} \nt{ListaMetodosInterfaz} \alt{} \eps}
\end{producciones}

\noindent{\small\color{grisMedio}La visibilidad se factoriza al frente de los tres tipos de miembro. Eso hace que el \texttt{public} del \nt{Constructor} original quede absorbido por \nt{VisibilidadOpcional}, y que \nt{Constructor} pase a competir con \nt{Tipo} por el token idClase: \texttt{C()\{\}} (constructor) vs \texttt{C x;} (atributo de tipo C).\par}

\begin{producciones}
\regla{Miembro}{\nt{VisibilidadOpcional} \nt{MiembroSinVisibilidad}}\hfill\mbox{\footnotesize\color{grisMedio}[+\,\texttt{visibilidad}]}
\regla{Miembro}{\nt{MiembroSinVisibilidad}}\hfill\mbox{\footnotesize\color{grisMedio}[-\,\texttt{visibilidad}]}
\regla{Miembro}{\nt{Constructor}}\hfill\mbox{\footnotesize\color{grisMedio}[-\,\texttt{visibilidad}]}
\regla{VisibilidadOpcional}{\tl{public} \alt{} \tl{private} \alt{} \eps}\hfill\mbox{\footnotesize\color{grisMedio}[+\,\texttt{visibilidad}]}
\end{producciones}

\noindent{\small\color{grisMedio}Sin el logro, el \texttt{public} sólo aparece en el constructor y es lo que lo distingue del resto de los miembros: es la regla original de la cátedra.\par}

\begin{producciones}
\regla{Constructor}{\tl{public} \tl{idClase} \nt{CuerpoDeMetodo}}\hfill\mbox{\footnotesize\color{grisMedio}[-\,\texttt{visibilidad}]}
\regla{MiembroSinVisibilidad}{\tl{static} \nt{TipoMetodo} \tl{idMetVar} \nt{CuerpoDeMetodo}}
\regla{MiembroSinVisibilidad}{\tl{void} \tl{idMetVar} \nt{CuerpoDeMetodo}}
\regla{MiembroSinVisibilidad}{\nt{TipoPrimitivo} \nt{DimensionesOpcionales} \tl{idMetVar} \nt{AtributoOMetodo}}
\regla{MiembroSinVisibilidad}{\tl{idGen} \nt{DimensionesOpcionales} \tl{idMetVar} \nt{AtributoOMetodo}}
\regla{MiembroSinVisibilidad}{\tl{idClase} \nt{MiembroDesdeIdClase}}
\end{producciones}

\noindent{\small\color{grisMedio}El NT insertado que decide entre constructor y miembro con tipo de referencia. Con visibilidad el constructor ya no lleva \texttt{public} propio, así que compite con el atributo/método de tipo referencia por el token idClase.\par}

\begin{producciones}
\regla{MiembroDesdeIdClase}{\nt{CuerpoDeMetodo}}\hfill\mbox{\footnotesize\color{grisMedio}[+\,\texttt{visibilidad}]}
\regla{MiembroDesdeIdClase}{\nt{TipoGenericoOpcional} \nt{DimensionesOpcionales} \tl{idMetVar} \nt{AtributoOMetodo}}
\end{producciones}

\noindent{\small\color{grisMedio}El NT insertado que decide entre atributo y método.\par}

\begin{producciones}
\regla{AtributoOMetodo}{\nt{InicializacionDeAtributo} \tl{;}}
\regla{AtributoOMetodo}{\nt{CuerpoDeMetodo}}
\end{producciones}

\noindent{\small\color{grisMedio}NT dedicado: la inicialización de un ATRIBUTO es un logro distinto que la de una variable local, y tienen que poder prenderse por separado.\par}

\begin{producciones}
\regla{InicializacionDeAtributo}{\tl{=} \nt{ExpresionCompuesta}}\hfill\mbox{\footnotesize\color{grisMedio}[+\,\texttt{atrinit}]}
\regla{InicializacionDeAtributo}{\eps}
\end{producciones}

\noindent{\small\color{grisMedio}Se comparte a propósito: es la misma construcción en método, método static y constructor.\par}

\begin{producciones}
\regla{CuerpoDeMetodo}{\nt{ArgsFormales} \nt{Bloque}}
\end{producciones}

\noindent{\small\color{grisMedio}Los métodos de interfaz no llevan visibilidad: el logro habla de Métodos, Atributos y Constructores, que son los \nt{Miembro} de una clase.\par}

\begin{producciones}
\regla{MetodoInterfaz}{\nt{TipoMetodo} \tl{idMetVar} \nt{ArgsFormales} \tl{;}}
\regla{TipoMetodo}{\nt{Tipo} \alt{} \tl{void}}
\regla{Tipo}{\nt{TipoBase} \nt{DimensionesOpcionales}}
\regla{TipoBase}{\nt{TipoPrimitivo} \alt{} \nt{TipoReferencia} \alt{} \tl{idGen}}
\regla{DimensionesOpcionales}{\tl{[} \tl{]} \nt{DimensionesOpcionales} \alt{} \eps}
\regla{TipoReferencia}{\tl{idClase} \nt{TipoGenericoOpcional}}
\end{producciones}

\noindent{\small\color{grisMedio}[L:float] — \texttt{float} es token del léxico desde la Etapa 1 pero no estaba en la gramática; sin esta alternativa el logro de Etapa 1 queda inutilizable.\par}

\begin{producciones}
\regla{TipoPrimitivo}{\tl{boolean} \alt{} \tl{char} \alt{} \tl{int}}
\regla{TipoPrimitivo}{\tl{float}}\hfill\mbox{\footnotesize\color{grisMedio}[+\,\texttt{float}]}
\end{producciones}

\noindent{\small\color{grisMedio}En POSICIÓN DE TIPO no se admite el diamante: \texttt{List<> x;} es inválido.\par}

\begin{producciones}
\regla{TipoGenericoOpcional}{\tl{<} \nt{ArgumentoDeTipo} \tl{>} \alt{} \eps}
\end{producciones}

\noindent{\small\color{grisMedio}[L:generico] — la recursión por TipoGenericoOpcional es la que habilita el anidamiento \texttt{List<List<Persona>>}. Sin el logro: \texttt{ArgumentoDeTipo -> idGen | idClase}.\par}

\begin{producciones}
\regla{ArgumentoDeTipo}{\tl{idGen}}
\regla{ArgumentoDeTipo}{\tl{idClase} \nt{TipoGenericoOpcional}}\hfill\mbox{\footnotesize\color{grisMedio}[+\,\texttt{generico}]}
\regla{ArgumentoDeTipo}{\tl{idClase}}\hfill\mbox{\footnotesize\color{grisMedio}[-\,\texttt{generico}]}
\regla{ArgsFormales}{\tl{(} \nt{ListaArgsFormalesOpcional} \tl{)}}
\regla{ListaArgsFormalesOpcional}{\nt{ListaArgsFormales} \alt{} \eps}
\regla{ListaArgsFormales}{\nt{ArgFormal} \nt{RestoListaArgsFormales}}
\regla{RestoListaArgsFormales}{\tl{,} \nt{ArgFormal} \nt{RestoListaArgsFormales} \alt{} \eps}
\regla{ArgFormal}{\nt{Tipo} \tl{idMetVar}}
\end{producciones}

\noindent{\small\color{grisMedio}Con [L:varclasica] una declaración local puede empezar con idClase (\texttt{C x;}), y una sentencia-expresión también (\texttt{C.m();}). Se factoriza ese idClase y el desempate lo hace el token siguiente: \texttt{.} es expresión, cualquier otro es declaración. La rama de expresión reusa \nt{ColaDeExpresionDesdePrimario} para no duplicar la jerarquía de expresiones.\par}

\begin{producciones}
\regla{Bloque}{\tl{\{} \nt{ListaSentencias} \tl{\}}}
\regla{ListaSentencias}{\nt{Sentencia} \nt{ListaSentencias} \alt{} \eps}
\regla{Sentencia}{\tl{;}}
\regla{Sentencia}{\nt{VarLocal} \tl{;}}
\regla{Sentencia}{\nt{Return} \tl{;}}
\regla{Sentencia}{\nt{If}}
\regla{Sentencia}{\nt{While}}
\regla{Sentencia}{\nt{Bloque}}
\regla{Sentencia}{\nt{For}}\hfill\mbox{\footnotesize\color{grisMedio}[+\,\texttt{for}]}
\end{producciones}

\noindent{\small\color{grisMedio}Sin variables locales clásicas, una sentencia-expresión puede empezar con idClase sin competir con nada: es la regla original.\par}

\begin{producciones}
\regla{Sentencia}{\nt{SentenciaExpresion} \tl{;}}\hfill\mbox{\footnotesize\color{grisMedio}[-\,\texttt{varclasica}]}
\regla{SentenciaExpresion}{\nt{Expresion}}\hfill\mbox{\footnotesize\color{grisMedio}[-\,\texttt{varclasica}]}
\end{producciones}

\noindent{\small\color{grisMedio}Con el logro, \texttt{C x;} y \texttt{C.m();} compiten por idClase y hay que factorizarlo.\par}

\begin{producciones}
\regla{Sentencia}{\nt{ExpresionSinIdClase} \tl{;}}\hfill\mbox{\footnotesize\color{grisMedio}[+\,\texttt{varclasica}]}
\regla{Sentencia}{\nt{TipoPrimitivo} \nt{DimensionesOpcionales} \nt{ListaDeclaradores} \tl{;}}\hfill\mbox{\footnotesize\color{grisMedio}[+\,\texttt{varclasica}]}
\regla{Sentencia}{\tl{idGen} \nt{DimensionesOpcionales} \nt{ListaDeclaradores} \tl{;}}\hfill\mbox{\footnotesize\color{grisMedio}[+\,\texttt{varclasica}]}
\regla{Sentencia}{\tl{idClase} \nt{SentenciaDesdeIdClase}}\hfill\mbox{\footnotesize\color{grisMedio}[+\,\texttt{varclasica}]}
\regla{SentenciaDesdeIdClase}{\nt{TipoGenericoOpcional} \nt{DimensionesOpcionales} \nt{ListaDeclaradores} \tl{;}}\hfill\mbox{\footnotesize\color{grisMedio}[+\,\texttt{varclasica}]}
\regla{SentenciaDesdeIdClase}{\tl{.} \tl{idMetVar} \nt{ArgsActuales} \nt{ColaDeExpresionDesdePrimario} \tl{;}}\hfill\mbox{\footnotesize\color{grisMedio}[+\,\texttt{varclasica}]}
\regla{VarLocal}{\tl{var} \tl{idMetVar} \tl{=} \nt{ExpresionCompuesta}}
\end{producciones}

\noindent{\small\color{grisMedio}Varias variables por sentencia, cada una con inicialización opcional: \texttt{int x,y,z = 10;}\par}

\begin{producciones}
\regla{ListaDeclaradores}{\tl{idMetVar} \nt{InicializacionDeDeclarador} \nt{RestoListaDeclaradores}}\hfill\mbox{\footnotesize\color{grisMedio}[+\,\texttt{varclasica}]}
\regla{RestoListaDeclaradores}{\tl{,} \tl{idMetVar} \nt{InicializacionDeDeclarador} \nt{RestoListaDeclaradores}}\hfill\mbox{\footnotesize\color{grisMedio}[+\,\texttt{varclasica}]}
\regla{RestoListaDeclaradores}{\eps}\hfill\mbox{\footnotesize\color{grisMedio}[+\,\texttt{varclasica}]}
\regla{InicializacionDeDeclarador}{\tl{=} \nt{ExpresionCompuesta} \alt{} \eps}\hfill\mbox{\footnotesize\color{grisMedio}[+\,\texttt{varclasica}]}
\regla{Return}{\tl{return} \nt{ExpresionOpcional}}
\regla{ExpresionOpcional}{\nt{Expresion} \alt{} \eps}
\regla{If}{\tl{if} \tl{(} \nt{Expresion} \tl{)} \nt{Sentencia} \nt{ElseOpcional}}
\regla{ElseOpcional}{\tl{else} \nt{Sentencia} \alt{} \eps}
\regla{While}{\tl{while} \tl{(} \nt{Expresion} \tl{)} \nt{Sentencia}}
\end{producciones}

\noindent{\small\color{grisMedio}[L:for] — Fors! E2 Una sola sentencia/expresión por separador. El \texttt{:} distingue el for de iteradores del estándar, y aparece justo después del declarador, así que el desempate es de un solo token.\par}

\begin{producciones}
\regla{For}{\tl{for} \tl{(} \nt{EncabezadoFor} \tl{)} \nt{Sentencia}}\hfill\mbox{\footnotesize\color{grisMedio}[+\,\texttt{for}]}
\regla{EncabezadoFor}{\tl{;} \nt{ExpresionOpcional} \tl{;} \nt{ExpresionOpcional}}\hfill\mbox{\footnotesize\color{grisMedio}[+\,\texttt{for}]}
\regla{EncabezadoFor}{\tl{var} \tl{idMetVar} \nt{EncabezadoForTrasVar}}\hfill\mbox{\footnotesize\color{grisMedio}[+\,\texttt{for}]}
\regla{EncabezadoFor}{\nt{TipoPrimitivo} \nt{DimensionesOpcionales} \tl{idMetVar} \nt{EncabezadoForTrasDeclarador}}\hfill\mbox{\footnotesize\color{grisMedio}[+\,\texttt{for}]}
\regla{EncabezadoFor}{\tl{idGen} \nt{DimensionesOpcionales} \tl{idMetVar} \nt{EncabezadoForTrasDeclarador}}\hfill\mbox{\footnotesize\color{grisMedio}[+\,\texttt{for}]}
\regla{EncabezadoFor}{\tl{idClase} \nt{EncabezadoForDesdeIdClase}}\hfill\mbox{\footnotesize\color{grisMedio}[+\,\texttt{for}]}
\regla{EncabezadoFor}{\nt{ExpresionSinIdClase} \nt{ClausulasFinalesFor}}\hfill\mbox{\footnotesize\color{grisMedio}[+\,\texttt{for}]}
\regla{EncabezadoForDesdeIdClase}{\nt{TipoGenericoOpcional} \nt{DimensionesOpcionales} \tl{idMetVar} \nt{EncabezadoForTrasDeclarador}}\hfill\mbox{\footnotesize\color{grisMedio}[+\,\texttt{for}]}
\regla{EncabezadoForDesdeIdClase}{\tl{.} \tl{idMetVar} \nt{ArgsActuales} \nt{ColaDeExpresionDesdePrimario} \nt{ClausulasFinalesFor}}\hfill\mbox{\footnotesize\color{grisMedio}[+\,\texttt{for}]}
\end{producciones}

\noindent{\small\color{grisMedio}El \texttt{:} distingue el for de iteradores del estándar, y está a un solo token.\par}

\begin{producciones}
\regla{EncabezadoForTrasVar}{\tl{:} \nt{Expresion}}\hfill\mbox{\footnotesize\color{grisMedio}[+\,\texttt{for}]}
\regla{EncabezadoForTrasVar}{\tl{=} \nt{ExpresionCompuesta} \nt{ClausulasFinalesFor}}\hfill\mbox{\footnotesize\color{grisMedio}[+\,\texttt{for}]}
\regla{EncabezadoForTrasDeclarador}{\tl{:} \nt{Expresion}}\hfill\mbox{\footnotesize\color{grisMedio}[+\,\texttt{for}]}
\regla{EncabezadoForTrasDeclarador}{\nt{InicializacionDeDeclaradorFor} \nt{ClausulasFinalesFor}}\hfill\mbox{\footnotesize\color{grisMedio}[+\,\texttt{for}]}
\regla{InicializacionDeDeclaradorFor}{\tl{=} \nt{ExpresionCompuesta} \alt{} \eps}\hfill\mbox{\footnotesize\color{grisMedio}[+\,\texttt{for}]}
\end{producciones}

\noindent{\small\color{grisMedio}Una sola expresión por separador, como pide el logro.\par}

\begin{producciones}
\regla{ClausulasFinalesFor}{\tl{;} \nt{ExpresionOpcional} \tl{;} \nt{ExpresionOpcional}}\hfill\mbox{\footnotesize\color{grisMedio}[+\,\texttt{for}]}
\end{producciones}

\noindent{\small\color{grisMedio}Capas, de menor a mayor precedencia: Expresion            asignación (NO asociativa, como en la gramática original) ExpresionCondicional ternario [L:ternario] ExpresionCompuesta   operadores binarios (un solo nivel, como el original) ExpresionBasica      operadores unarios prefijos Operando             literal, o referencia con posfijo opcional [L:posfijo]\par}

\begin{producciones}
\regla{Expresion}{\nt{ExpresionCondicional} \nt{RestoExpresion}}
\end{producciones}

\noindent{\small\color{grisMedio}El lado derecho es ExpresionCondicional y NO Expresion: la gramática original pone \nt{ExpresionCompuesta} ahí, o sea que \texttt{a = b = c} NO es válido en MiniJava.\par}

\begin{producciones}
\regla{RestoExpresion}{\nt{OperadorAsignacion} \nt{ExpresionCondicional} \alt{} \eps}
\regla{OperadorAsignacion}{\tl{=}}
\regla{ExpresionCondicional}{\nt{ExpresionCompuesta} \nt{RestoCondicional}}
\end{producciones}

\noindent{\small\color{grisMedio}Con el logro apagado \nt{RestoCondicional} queda en solo ε, o sea que \nt{ExpresionCondicional} se vuelve un alias de \nt{ExpresionCompuesta} y la capa desaparece sin dejar rastro. Asociativo a derecha, como en Java.\par}

\begin{producciones}
\regla{RestoCondicional}{\tl{?} \nt{ExpresionCondicional} \tl{:} \nt{ExpresionCondicional}}\hfill\mbox{\footnotesize\color{grisMedio}[+\,\texttt{ternario}]}
\regla{RestoCondicional}{\eps}
\regla{ExpresionCompuesta}{\nt{ExpresionBasica} \nt{RestoExpresionCompuesta}}
\regla{RestoExpresionCompuesta}{\nt{OperadorBinario} \nt{ExpresionBasica} \nt{RestoExpresionCompuesta} \alt{} \eps}
\regla{OperadorBinario}{\tl{||} \alt{} \tl{\&\&} \alt{} \tl{==} \alt{} \tl{!=} \alt{} \tl{<} \alt{} \tl{>} \alt{} \tl{<=} \alt{} \tl{>=} \alt{} \tl{+} \alt{} \tl{-} \alt{} \tl{*} \alt{} \tl{/} \alt{} \tl{\%}}
\regla{ExpresionBasica}{\nt{OperadorUnario} \nt{Operando}}
\regla{ExpresionBasica}{\nt{Operando}}
\regla{OperadorUnario}{\tl{+} \alt{} \tl{-} \alt{} \tl{!}}
\regla{Operando}{\nt{Primitivo}}
\regla{Operando}{\nt{Referencia} \nt{PosfijoOpcional}}
\end{producciones}

\noindent{\small\color{grisMedio}Va sobre \nt{Referencia} y no sobre \nt{Operando}: \texttt{5++} no tiene sentido.\par}

\begin{producciones}
\regla{PosfijoOpcional}{\tl{++}}\hfill\mbox{\footnotesize\color{grisMedio}[+\,\texttt{posfijo}]}
\regla{PosfijoOpcional}{\tl{-{}-}}\hfill\mbox{\footnotesize\color{grisMedio}[+\,\texttt{posfijo}]}
\regla{PosfijoOpcional}{\eps}
\end{producciones}

\noindent{\small\color{grisMedio}[L:float] — floatLiteral\par}

\begin{producciones}
\regla{Primitivo}{\tl{true} \alt{} \tl{false} \alt{} \tl{intLiteral} \alt{} \tl{charLiteral} \alt{} \tl{null}}
\regla{Primitivo}{\tl{floatLiteral}}\hfill\mbox{\footnotesize\color{grisMedio}[+\,\texttt{float}]}
\end{producciones}

\noindent{\small\color{grisMedio}--- Expresión que no puede empezar con idClase ---------------------------- Se usa donde una expresión compite con una declaración de variable local (sentencia y encabezado del for). Sólo duplica las *cabezas*: todas las colas se comparten con la jerarquía normal, así que no hay lógica repetida.\par}

\begin{producciones}
\regla{ExpresionSinIdClase}{\nt{OperadorUnario} \nt{Operando} \nt{ColaTrasOperando}}
\regla{ExpresionSinIdClase}{\nt{Primitivo} \nt{ColaTrasOperando}}
\regla{ExpresionSinIdClase}{\nt{PrimarioSinIdClase} \nt{ColaDeExpresionDesdePrimario}}
\regla{ColaTrasOperando}{\nt{RestoExpresionCompuesta} \nt{RestoCondicional} \nt{RestoExpresion}}
\regla{ColaDeExpresionDesdePrimario}{\nt{RestoReferencia} \nt{PosfijoOpcional} \nt{ColaTrasOperando}}
\regla{Referencia}{\nt{Primario} \nt{RestoReferencia}}
\regla{RestoReferencia}{\nt{Encadenado} \nt{RestoReferencia}}
\regla{RestoReferencia}{\nt{AccesoArreglo} \nt{RestoReferencia}}
\regla{RestoReferencia}{\eps}
\regla{Encadenado}{\tl{.} \tl{idMetVar} \nt{EncadenadoAtributoOMetodo}}
\end{producciones}

\noindent{\small\color{grisMedio}Dedicado a este contexto: vacío = era <VarEncadenada>, con argumentos = era <MetodoEncadenado>.\par}

\begin{producciones}
\regla{EncadenadoAtributoOMetodo}{\nt{ArgsActuales} \alt{} \eps}
\regla{AccesoArreglo}{\tl{[} \nt{Expresion} \tl{]}}
\regla{Primario}{\nt{LlamadaMetodoEstatico}}
\regla{Primario}{\nt{PrimarioSinIdClase}}
\regla{LlamadaMetodoEstatico}{\tl{idClase} \tl{.} \tl{idMetVar} \nt{ArgsActuales}}
\regla{PrimarioSinIdClase}{\tl{this}}
\regla{PrimarioSinIdClase}{\tl{stringLiteral}}
\regla{PrimarioSinIdClase}{\nt{ParentizadaOLambda}}
\regla{PrimarioSinIdClase}{\nt{AccesoVarLlamadaOLambda}}
\regla{PrimarioSinIdClase}{\nt{Creacion}}
\end{producciones}

\noindent{\small\color{grisMedio}Dedicado a este contexto: vacío = era <AccesoVar>, con argumentos = era <LlamadaMetodo>, con flecha = lambda de un parámetro sin paréntesis.\par}

\begin{producciones}
\regla{AccesoVarLlamadaOLambda}{\tl{idMetVar} \nt{AccesoLlamadaOLambda}}
\regla{AccesoLlamadaOLambda}{\nt{ArgsActuales}}
\regla{AccesoLlamadaOLambda}{\tl{->} \nt{CuerpoLambda}}\hfill\mbox{\footnotesize\color{grisMedio}[+\,\texttt{lambda}]}
\regla{AccesoLlamadaOLambda}{\eps}
\end{producciones}

\noindent{\small\color{grisMedio}[L:lambda] — MiniLambdas E2 \texttt{(} abre tanto una expresión parentizada como la lista de parámetros de una lambda. Parece necesitar lookahead ilimitado (\texttt{(a,b,c) ->} vs \texttt{(a+b)}), pero NO lo necesita: - \texttt{( )}      sólo puede ser una lambda de cero parámetros; - \texttt{( e , …}  sólo puede ser una lambda: MiniJava no tiene operador coma; - \texttt{( e )}    se decide con el token que sigue al \texttt{)}: \texttt{->} es lambda, cualquier otro es una expresión parentizada. Con eso alcanza un token de lookahead en cada punto.\par}

\begin{producciones}
\regla{ParentizadaOLambda}{\tl{(} \nt{ContenidoDeParentesis}}
\regla{ContenidoDeParentesis}{\tl{)} \tl{->} \nt{CuerpoLambda}}\hfill\mbox{\footnotesize\color{grisMedio}[+\,\texttt{lambda}]}
\regla{ContenidoDeParentesis}{\nt{Expresion} \nt{RestoDeParentesis}}
\regla{RestoDeParentesis}{\tl{)} \nt{FinDeParentizadaOLambda}}
\regla{RestoDeParentesis}{\tl{,} \nt{ListaParametrosLambda} \tl{)} \tl{->} \nt{CuerpoLambda}}\hfill\mbox{\footnotesize\color{grisMedio}[+\,\texttt{lambda}]}
\regla{FinDeParentizadaOLambda}{\tl{->} \nt{CuerpoLambda}}\hfill\mbox{\footnotesize\color{grisMedio}[+\,\texttt{lambda}]}
\regla{FinDeParentizadaOLambda}{\eps}
\regla{ListaParametrosLambda}{\tl{idMetVar} \nt{RestoListaParametrosLambda}}\hfill\mbox{\footnotesize\color{grisMedio}[+\,\texttt{lambda}]}
\regla{RestoListaParametrosLambda}{\tl{,} \tl{idMetVar} \nt{RestoListaParametrosLambda} \alt{} \eps}\hfill\mbox{\footnotesize\color{grisMedio}[+\,\texttt{lambda}]}
\end{producciones}

\noindent{\small\color{grisMedio}Una única expresión, sin llaves.\par}

\begin{producciones}
\regla{CuerpoLambda}{\nt{Expresion}}\hfill\mbox{\footnotesize\color{grisMedio}[+\,\texttt{lambda}]}
\regla{Creacion}{\tl{new} \nt{ArregloOInstancia}}
\regla{ArregloOInstancia}{\nt{TipoPrimitivo} \nt{CreacionDeArreglo}}
\regla{ArregloOInstancia}{\tl{idGen} \nt{CreacionDeArreglo}}
\regla{ArregloOInstancia}{\tl{idClase} \nt{CreacionDesdeIdClase}}
\regla{CreacionDesdeIdClase}{\tl{<} \nt{GenericoDeCreacion}}
\regla{CreacionDesdeIdClase}{\nt{CreacionDeArreglo}}
\regla{CreacionDesdeIdClase}{\nt{ArgsActuales}}
\end{producciones}

\noindent{\small\color{grisMedio}[L:generico] — el diamante \texttt{<>} sólo se admite acá, en la rama que va a ArgsActuales: \texttt{new C<>()} es válido y \texttt{new C<>[3]} no. La restricción del logro queda impuesta por la gramática, no delegada al semántico.\par}

\begin{producciones}
\regla{GenericoDeCreacion}{\tl{>} \nt{ArgsActuales}}\hfill\mbox{\footnotesize\color{grisMedio}[+\,\texttt{generico}]}
\regla{GenericoDeCreacion}{\nt{ArgumentoDeTipo} \tl{>} \nt{ArregloOInstanciaDeClase}}
\regla{ArregloOInstanciaDeClase}{\nt{CreacionDeArreglo}}
\regla{ArregloOInstanciaDeClase}{\nt{ArgsActuales}}
\end{producciones}

\noindent{\small\color{grisMedio}[L:arrinit] — \texttt{new int[]\{1,2,3\}}. Sin el logro: CreacionDeArreglo -> [ Expresion ] DimensionesConTamanioOpcional\par}

\begin{producciones}
\regla{CreacionDeArreglo}{\tl{[} \nt{RestoCreacionDeArreglo}}
\regla{RestoCreacionDeArreglo}{\nt{Expresion} \tl{]} \nt{DimensionesConTamanioOpcional}}
\regla{RestoCreacionDeArreglo}{\tl{]} \tl{\{} \nt{ListaExpsOpcional} \tl{\}}}\hfill\mbox{\footnotesize\color{grisMedio}[+\,\texttt{arrinit}]}
\regla{DimensionesConTamanio}{\tl{[} \nt{Expresion} \tl{]} \nt{DimensionesConTamanioOpcional}}
\regla{DimensionesConTamanioOpcional}{\nt{DimensionesConTamanio} \alt{} \eps}
\regla{ArgsActuales}{\tl{(} \nt{ListaExpsOpcional} \tl{)}}
\regla{ListaExpsOpcional}{\nt{ListaExps} \alt{} \eps}
\regla{ListaExps}{\nt{Expresion} \nt{RestoListaExps}}
\regla{RestoListaExps}{\tl{,} \nt{Expresion} \nt{RestoListaExps} \alt{} \eps}
\end{producciones}



% ===========================================================================
\section{El analizador sintáctico}
\label{sec:analizador}
% ===========================================================================

\subsection{Diseño}

Descendente recursivo predictivo con el Esquema Simple de Clase 5: un método por no terminal,
\texttt{match} para los terminales, y la producción se elige mirando si el token actual está en los
Primeros de su cuerpo. Si ninguna coincide, el \texttt{else} final lanza el error; si alguna parte
derecha deriva a \eps, es un \texttt{else} vacío con esa producción citada. La recursión a derecha se
traduce con recursión.

\begin{center}\small
\begin{tabular}{@{}>{\raggedright\arraybackslash}p{0.40\linewidth}>{\raggedright\arraybackslash}p{0.54\linewidth}@{}}
\toprule
\textbf{Clase} & \textbf{Responsabilidad} \\
\midrule
\texttt{AnalizadorSintactico} & Carga el primer token, llama a \nt{Inicial} y devuelve los errores. \\
\texttt{EntradaDeTokens} & El token actual y \texttt{match}: lo único que consume entrada. \\
\texttt{Sintaxis} & Conecta las nueve secciones, que son mutuamente recursivas. \\
\texttt{partes/AnalizadorDe\{Clases, Miembros, Tipos, Sentencias, Fors, Expresiones, Referencias,
Creacion, Argumentos\}} & Un método por no terminal de cada sección de la gramática. \\
\texttt{primeros/PrimerosDeX} & Los Primeros de cada producción, generados desde la gramática. \\
\texttt{recuperacion/*} & Modo Pánico. \\
\texttt{errores/ExcepcionSintactica} & El error sintáctico. \\
\bottomrule
\end{tabular}
\end{center}

\subsection{Errores que detecta}
\label{sec:errores}

El Esquema Simple falla en dos lugares: un \texttt{match} que no coincide (se esperaba un terminal y
llegó otro) y un no terminal no anulable cuyo token actual no está en los Primeros de ninguna de sus
producciones (se esperaba alguno de esos terminales). El mensaje nombra lo esperado y lo encontrado:

\begin{consola}
Error Sintactico en linea 3: se esperaba idClase se encontro "clasesita"

[Error:clasesita|3]
\end{consola}

El lexema del fin de archivo es \tl{\$}. Con Modo Pánico, ante un error se descarta entrada hasta un
\tl{;}, una \tl{\}}, una \tl{\{} que no abre un elemento nuevo, o un token que empiece el siguiente
elemento de la lista (clases, miembros, métodos de interfaz o sentencias), y se reporta cada error
en orden. El primero es siempre el mismo que sin recuperación. Los errores léxicos los reporta el
analizador de la Etapa 1.

% ===========================================================================
\section{Decisiones}
\label{sec:decisiones}
% ===========================================================================

\subsection{Sobre puntos consultados}

{\small
\begin{longtable}{@{}>{\raggedright\arraybackslash}p{0.30\linewidth}>{\raggedright\arraybackslash}p{0.64\linewidth}@{}}
\toprule
\textbf{Tema} & \textbf{Decisión} \\
\midrule\endhead
Huecos del \tl{for} & Una expresión en cada uno, obligatoria: \texttt{for(;;)} y
\texttt{for(var i = 0; i < n; )} no son válidos. \\
Variable de control del \tl{for} & Sólo con \tl{var}: \texttt{for (int i = 0; \dots)} y
\texttt{for (Persona p : lista)} no son válidos. \\
Asignaciones anidadas & El lado derecho de un \tl{=}, en la asignación y en los inicializadores, es
una \nt{ExpresionCondicional}: \texttt{x = y = 2} e \texttt{int x = y = 2;} no son válidos;
\texttt{x = (y = 2)} sí. \\
Inicialización de arreglos & Como en Java: con inicializador ninguna dimensión lleva tamaño, se
anidan y admiten coma final. Sin inicializador rige la base, con toda dimensión con tamaño:
\texttt{new int[2][]} no es válido. \\
Atributos & \tl{static} sólo modifica métodos, y un atributo por declaración: \texttt{static int x;}
e \texttt{int a, b;} no son válidos. \\
Modificadores & \tl{static} y la visibilidad en cualquier orden, sin repetir; \tl{protected} no
existe. \\
Logros & No son configurables. \\
Recursión a derecha & Se traduce con recursión, no con \texttt{while}. \\
\texttt{else} final & Ningún método de no terminal lo omite. \\
Lexema del EOF & \tl{\$}. \\
\bottomrule
\end{longtable}}

\subsection{Otras}

{\small
\begin{longtable}{@{}>{\raggedright\arraybackslash}p{0.30\linewidth}>{\raggedright\arraybackslash}p{0.64\linewidth}@{}}
\toprule
\textbf{Tema} & \textbf{Decisión} \\
\midrule\endhead
Métodos de interfaz & Admiten \tl{public} y no \tl{private}: Visibilidad Mejorada nombra a los
métodos, y en Java un método de interfaz privado exige cuerpo, que \nt{MetodoInterfaz} no tiene. \\
Operadores posfijos & Sobre una \nt{Referencia} y uno solo: \texttt{5++} y \texttt{x++ ++} no son
válidos. Que la referencia sea una variable lo verifica el semántico. \\
\tl{float} & \tl{float} y \tl{floatLiteral} entran a \nt{TipoPrimitivo} y \nt{Primitivo}, para que
el logro Floats I tenga uso en la sintaxis. \\
Ternario & Operandos del JLS §15.25 y asociativo a derecha: \texttt{a ? b = 1 : c} es válido y
\texttt{a ? b : c = d} no. \\
\bottomrule
\end{longtable}}

\end{document}
